\documentclass[10pt,letterpaper]{article}
\usepackage[spanish]{babel}
\usepackage{listings}

% Preámbulo

\title{Sistema de gestión de una Biblioteca Universitaria}
\author{Shomara Berenice Rivera Huerta}
\date{\today}

\usepackage{xcolor}

\definecolor{codegreen}{rgb}{0,0.6,0}
\definecolor{codegray}{rgb}{0.5,0.5,0.5}
\definecolor{codepurple}{rgb}{0.58,0,0.82}
\definecolor{backcolour}{rgb}{0.95,0.95,0.92}

\lstdefinestyle{mystyle}{
    backgroundcolor=\color{backcolour},   
    commentstyle=\color{codegreen},
    keywordstyle=\color{magenta},
    numberstyle=\tiny\color{codegray},
    stringstyle=\color{codepurple},
    basicstyle=\ttfamily\footnotesize,
    breakatwhitespace=false,         
    breaklines=true,                 
    captionpos=b,                    
    keepspaces=true,                 
    numbers=left,                    
    numbersep=5pt,                  
    showspaces=false,                
    showstringspaces=false,
    showtabs=false,                  
    tabsize=2
}
\lstset{style=mystyle}

 \renewcommand{\lstlistingname}{Código}
 \renewcommand{\lstlistlistingname}{Lista de códigos}

\begin{document}
\maketitle
\tableofcontents
\lstlistoflistings

\section{Introducción}

La biblioteca cuenta con \textbf{libro}, de los cuales puede tener varios ejemplares físicos para cada título. Cada libro tienen un título, año de publicación, ISBN y pertenece a una o mas categorías (Ciencias, Humanidades, Ingeniería, etc.). Además, cada libro fue escrito por uno o más autores.

Los usuarios pueden ser estudiantes o profesores. De cada uno se guarda el nombre, correo , número de cuenta y tipo. Un usuario puede tener activos varios préstamos simultáneos (máximo 3), pero un ejemplar solo puede estar prestado a un usuario a la vez.

Cada préstamo registra la fecha de salida, la fecha de devolución esperada y la fecha real de devolución (en el caso de que ya se haya devuelto el ejemplar). Si un préstamo se devuelve más tarde, se registra una multa con el monto y si fue pagada o no.

\section{Descripción del problema}
\subsection{Análisis y diseño}
\subsection{Implementación y validación}


% Incluir los scripts solo si existen para evitar errores en la compilación
\IfFileExists{codes/create.sql}{%
    \lstinputlisting[language=SQL, frame=single, numbers=left, caption={Script de creación de las tablas para el modelo de datos}]{codes/create.sql}%
}{%
    \begin{center}\textbf{Archivo codes/create.sql no encontrado.}\end{center}%
}
\IfFileExists{codes/insertions.sql}{%
    \lstinputlisting[language=SQL, frame=single, numbers=left, caption={Script de inserción de datos sobre la estructura del modelo de datos}]{codes/insertions.sql}%
}{%
    \begin{center}\textbf{Archivo codes/insertions.sql no encontrado.}\end{center}%
}

\section{Conclusiones}

\section{Bibliografía}


\end{document}